% Separate propositions sheet. Replace all prompts with defensible propositions.
\documentclass[12pt,a4paper]{article}
\usepackage{iftex}
\ifPDFTeX\usepackage[T1]{fontenc}\usepackage{lmodern}\else
\usepackage{fontspec}\setmainfont{Latin Modern Roman}\fi
\usepackage[margin=28mm]{geometry}
\usepackage[british]{babel}
\usepackage{enumitem}
\usepackage[hidelinks,pdflang=en-GB]{hyperref}
% Edit these fields first. Keep personal research in a private project.
\newcommand{\ThesisTitle}{Law, authority and accountable institutions}
\newcommand{\ThesisSubtitle}{A working template for doctoral legal research}
\newcommand{\ThesisAuthor}{Your name}
\newcommand{\ThesisUniversity}{Your university}
\newcommand{\ThesisFaculty}{Faculty of Law}
\newcommand{\ThesisDegree}{Doctoral thesis working example}
\newcommand{\ThesisDate}{8 November 2026}
\newcommand{\ThesisSupervisors}{Supervisor names and approved titles}
\newcommand{\ThesisLegalCutoff}{8 November 2026 (demonstration date)}
\newcommand{\ThesisKeywords}{law, doctoral research, legal methodology}
\newcommand{\ThesisDocumentLabel}{LEXTHESIS / WORKING EXAMPLE}
\newcommand{\ThesisTranslatedTitle}{Recht, gezag en verantwoordelijke instituties}
\newcommand{\ThesisTranslatedSubtitle}{Een werkmodel voor rechtswetenschappelijk promotieonderzoek}

\hypersetup{pdftitle={Propositions accompanying \ThesisTitle},pdfauthor={\ThesisAuthor}}
\begin{document}
\thispagestyle{empty}
\begin{center}
{\Large\bfseries Propositions}\par\medskip
accompanying\par\medskip
{\large\ThesisTitle}\par\medskip
\ThesisAuthor
\end{center}
\bigskip
\noindent\textbf{Working sheet.} Replace the six prompts below with academically
defensible propositions and obtain any required approval. These prompts are not
completed propositions. Groningen's March 2026 regulations require at least six
propositions accompanying the thesis on a separate sheet.
\bigskip
\begin{enumerate}[itemsep=12pt]
\item State a precise principal finding established by your research.
\item State a second finding and identify the distinction on which it depends.
\item Formulate a defensible implication of your interpretation of the law.
\item State a proposition addressing a significant counterargument.
\item Formulate a methodological insight supported by your research experience.
\item State a further academically defensible proposition within the scope allowed by your university.
\end{enumerate}
\vfill
\noindent\small Verify the applicable rules, wording and approval procedure
with your supervisor and Graduate School before submission.
\end{document}
